% No preamble. Follows cv- concrete inputs and cl- support localization.
% Unique prefix ci-.
\section{One integration of the classical action at every earlier depth}
\label{ci-section}

The support theorem upgrades the bare calculation of Notebook 20.
Its exact scope matters: $S_M$ is the classical action obtained
by minimizing the original Wilson action through the earlier
ratio $M$. The following integral uses $S_M$ and the source
gauge and measure factors. It has no previously generated
quantum correction $E_k$. Thus the result is uniform through
classical depth, not a claim that the quantum induction is
already closed.

Let $\Phi_V$ be Notebook 19's exact retained chart, let
$a=\log V/i$, and put
\begin{equation}\label{ci-action}
 \mathcal T_{M,V}(x)=
 S_M\bigl(\log(e^{i\Phi_V(x)}V)/i\bigr)
                    +\mathcal G(e^{i\Phi_V(x)}V).
\end{equation}
The full product-domain Hessian is
$A_M(V)=D_x^2\mathcal T_{M,V}(0)$, including the absolute
gauge function away from its zero slice. Its flat baseline is
$A_{M,0}=K_L^T(q_M+T_L^TT_L)K_L$.
The preceding support construction gives local scalar
coefficients $\mathcal T_{M,Y}$ and symmetric precision
coefficients with all their retained rows, columns and
background dependence in $Y$.

Group unit current cells into fixed $L$-blocks, assigning their
retained positive bonds to them. The Gaussian cell dimension
is then $D=d(L^d-1)\dim\mathfrak g$, independent of $M$ and
volume. Any bounded change of support convention and anchor
multiplicity in this grouping costs constants depending only
on $L$. After the spatial-exponent and cardinality reserves
proved above, the resulting coefficients have exactly the
weighted support norms of Notebook 18.

Define the logarithmic interaction $v_{M,Y}$ by subtracting
the constant, linear and quadratic terms from each
$-g^{-2}\mathcal T_{M,Y}(a,gz)$ and adding the normalized
straightening-Jacobian and Haar logarithms of Notebook 19.
The gauge Taylor remainder is already included in
$\mathcal T_{M,Y}$ and is not counted twice. There is a
common small product background ball such that, for every
fixed required support-cardinality weight and a sufficiently
small spatial exponent,
\begin{align}
 \nu_M(b)&\le C_L b,\label{ci-inputs}\\
 \sup_i\sum_{Y\ni i}e^{a_0|Y|+\kappa\ell(Y)}
                      \|v_{M,Y}\|_{b,R}
       &\le C_{L,a_0,\kappa}gR^3,
       \qquad R\ge1,\quad gR<r_L.\notag
\end{align}
Here $\nu_M$ is the precision support norm, not just the
norm of the summed matrix, and $\ell$ is the metric tree
length used by Notebook 18. All constants are uniform in
the earlier classical depth and periodic volume.

\begin{theorem}[Uniform classical-reference integration]
\label{ci-theorem}
Fix odd $L$, the gauge group, $c>0$ and $p>1/2$, and set
$R(g)=c(\log g^{-2})^p$. There are
$b,g_0,a_*,C,c_*>0$, independent of $M=L^k$ and compatible
periodic volume, for which
\begin{equation}\label{ci-normalized}
 Z_{M,g}(V)=\int d\mu_{A_M(V)^{-1}}(z)\,
       \chi_{R(g)}(z)\exp\left(\sum_Yv_{M,Y}(V,z)\right)
\end{equation}
has a holomorphic logarithm on the independent current
background product ball $\|a\|_\infty<b$, for $0<g<g_0$.
There is a convergent connected-support expansion
\begin{equation}\label{ci-result}
 \log Z_{M,g}(V)=\sum_HE_{M,g,H}(V|_H),\qquad
 \sup_i\sum_{H\ni i}e^{a_*|H|}\|E_{M,g,H}\|_b
       \le C\{gR(g)^3+e^{-c_*R(g)^2}\}.
\end{equation}
Its canonical difference from the pure Gaussian cutoff
expansion has norm at most $CgR(g)^3$.
\end{theorem}
\begin{proof}
Notebook 19 supplies the flat spectral constants and weighted
inverse uniformly in $M$. Choose the spatial exponents in
Notebook 18 smaller than both that inverse exponent and half
the available Steiner exponent for the support coefficients.
Choose a fixed $0<\theta\le1/4$, for example $1/(4D)$.
The arbitrary fixed cardinality reserve in the support
theorem pays $4+\log\mathcal K_\theta$ in its interaction norm.
The first bound in \eqref{ci-inputs} permits one
$M$-independent $b>0$ satisfying every local-precision
smallness condition of its tunable strictly local theorem.

The interaction parameter is at most $CgR^3$ by the second
bound in \eqref{ci-inputs}. The bad-cell parameter is at most
$Ce^{-c_*R^2}$ by the correlated real-baseline comparison.
The retained-bond rectangles are covered by the cell tail
argument of Notebook 20 and remain fixed under the exact
chart. A common $g_0$ makes both quarter budgets hold and
keeps the recovered whole field within its analytic radius.
Apply Notebook 18 and its periodic extension.
\end{proof}

At a physical background $V_*(W)$, the sum of the separated
linear terms is exactly zero and the gauge constant is zero.
Thus \eqref{ci-normalized} is the actual fluctuation factor
of the constrained classical-reference integral in that
chart. On the larger product domain it is the explicitly
centered analytic extension, as in Notebook 20.

The normalization factors are equally controlled. With
$B_M(V)=A_M(V)-A_{M,0}$ and $C_{M,0}=A_{M,0}^{-1}$, set
\begin{equation}\label{ci-det}
 \mathcal Q_M(V)=-\tfrac12\log\det(A_M(V)A_{M,0}^{-1})
   =\tfrac12\sum_{n\ge1}\frac{(-1)^n}{n}
                         \operatorname{Tr}[(C_{M,0}B_M(V))^n].
\end{equation}
The closed-chain proof of Notebook 20 uses only weighted
covariance row sums and the precision-label norm.
Equation \eqref{ci-inputs} supplies those uniformly in $M$,
even though the precision supports are no longer bounded
in size. Joining each label by a minimal tree and each
covariance link by a shortest path gives the same bound.
Consequently $\mathcal Q_M$ has strictly local connected
activities with anchored norm at most $Cb$. The separate
delta factor
$\log[\mathcal D(V)/\mathcal D(1)]=-\sum_c\operatorname{Tr}
\log\pi_c(V)$ has fixed-range anchored norm at most $Cb$.
Neither normalization term is silently counted as $O(g)$.

For real physical backgrounds define $I^{\rm sf}_{M,g}(V)$
to be the constrained current group integral with weight
$\exp[-g^{-2}(S_M(\log U/i)+\mathcal G(U))]$, restricted
to the local retained chart and its specified rectangle.
Using the notation $N_b,N_c,n,c_H,\mathcal D$ from
Notebook 20, its exact formula is
\begin{equation}\label{ci-exact}
 \begin{split}
 I^{\rm sf}_{M,g}(V)={}&
 c_H^{N_b-N_c}g^n\mathcal D(1)
       (2\pi)^{n/2}(\det A_{M,0})^{-1/2}
                       e^{-S_M(a)/g^2}\\
 &\times\exp\left\{
 \log\frac{\mathcal D(V)}{\mathcal D(1)}
       +\mathcal Q_M(V)+\sum_HE_{M,g,H}(V|_H)\right\}.
 \end{split}
\end{equation}
There is no unaccounted field-dependent normalization in
this defined integral. A gauge-insertion identity outside
the chart would additionally require $S_M$ itself to be
defined and gauge invariant on the chosen admissible
rooted-gauge-invariant domain. The small-potential construction
does not assert that extension. The exact chart formula above
does not require it.

\section{The next induction problem}
\label{ci-next}

The classical-reference integral is now controlled uniformly
over the number of earlier classical averaging steps.
This removes a separate localization premise from the
homogeneous small-field argument. Its support statement
initially records dependence on independent current potentials.
Notebook 22 supplies the further strict support expansion of
the actual stationary background $V_*(W)$ and transports the
generated activities into $W$, with explicit radius and
support-weight losses. Notebook 23 proves local gauge
invariance of the complete correction, including the
normalization factors, and its quadratic Ward identity.

Notebook 24 then supplies a conditional inherited-interaction
bound, and Notebook 26 meets its support-weight test by using
a persistent auxiliary-cell convention. The next problem is
stable analytic domains and control of the accumulated
interaction sizes, with extraction of vacuum and marginal
terms and estimates along the required coupling trajectory.
The bounded one-step estimates do not supply this recurrence.
Notebook 25's smaller representative gives one analytic
continuation tool for addressing the domain problem.
Large-field complements and mixed histories must be included
in the same induction. The continuum observable bounds,
symmetry restoration and reconstruction remain the final
existence obligations recorded in Notebook 01.
